\section{Introduction}
Low-altitude economy (LAE) services use UAVs for sensing, delivery, monitoring, and emergency support. In such services, UAV swarms can act as flying edge nodes for IoT users and vehicles. They can relay data, serve local requests, cache popular content, and provide temporary communication coverage. UAV-assisted wireless networks are also a key topic in 5G and beyond systems because UAVs provide flexible deployment and good line-of-sight opportunities \cite{li2019uavcommunications}.

The main difficulty is that the service state changes quickly. Users move, content requests are bursty, UAV energy is limited, and air-to-ground links are affected by distance, altitude, speed mismatch, and weather. A scheduler that only maximizes the current rate may switch swarms too often. A scheduler that only minimizes energy may delay urgent demand. A cache-aware method may still fail when the cached content cannot be served over a risky link. Hence, computation, communication, and caching, called 3C resources in this paper, should be handled together.

Existing works address important parts of this problem. MUCCO jointly considers coded caching, user grouping, and UAV deployment for energy-efficient data delivery \cite{wei2025mucco}. A recent post-disaster UAV-swarm work uses collaborative beamforming, multi-path routing, and rate maximization for emergency communication \cite{zheng2025postdisaster}. Dynamic task collaboration studies asynchronous two-layer UAV-swarm control and stability for task-volume consensus \cite{zhou2026dynamic}. These works are useful references, but they do not directly address predictive cache-value-aware 3C scheduling under environmental link uncertainty and competing dynamic content demand.

This paper proposes \pchorch, a predictive 3C orchestration framework for heterogeneous UAV swarms. The framework uses an ANN to predict next-slot link margin and converts it into outage risk using weather-wise calibration. It then scores feasible swarm-cluster assignments using predicted link risk, useful cache value, delay, energy, load, fairness, and mobility-aware dwell cost. The main practical method is \pclr. We also study \psr, an internal stability-regularized variant, and \etp, an exploratory event-triggered ablation.

The contributions are as follows.
\begin{itemize}
    \item We build a predictive 3C orchestration model for heterogeneous UAV swarms serving moving low-altitude IoT demand under clear, rainy, and rain-hot weather.
    \item We integrate ANN-based link-risk prediction with weather-wise uncertainty calibration and use it inside the scheduling decision.
    \item We design \pclr, a cache-value-aware and mobility-aware controller that jointly considers link margin, outage, energy, delay, load, and handover cost.
    \item We report \psr\ as a stability-regularized internal variant and \etp\ as an exploratory ablation, without hiding negative or mixed results.
    \item We evaluate all methods on 80 blind paper-lock seeds across seven regimes and report confidence intervals and Holm-corrected statistical tests.
\end{itemize}

The rest of the paper is organized as follows. Section II reviews related work. Section III gives the system model. Section IV describes the ANN link-risk predictor. Sections V and VI present \pchorch\ and its algorithms. Section VII explains the experimental setup. Section VIII discusses the results. Section IX gives ablations and limitations. Section X concludes the paper.
